% main.tex
% Universitaire Papers - Full Documentation & Usage Guide
% Author: Matt
% Purpose: A self-contained manual that explains paperSettings.tex,
%          standardCommands.tex and demonstrates how to use the macros.
%
% Save this file as main.tex in the root of the repo alongside:
%   - paperSettings.tex
%   - standardCommands.tex
%   - main.bib
%
% Compile with: buildLatexDocument.sh (or fullBuild.sh)
% Or upload to Overleaf and compile there.
% -----------------------------------------

% ----------- Document settings -----------
\documentclass[nonacm, sigconf, balance=truu]{acmart}

%----------- DONT TOUCH BELOW THIS LINE -----------
\usepackage{xparse}
\usepackage{xstring}

% ----------- Encoding & Fonts -----------
\usepackage[T1]{fontenc}
\usepackage[utf8]{inputenc}
\usepackage{scalefnt}

% ----------- Page layout -----------
\usepackage{expl3}
\usepackage{array}
\setlength{\columnsep}{0.333in}
\renewcommand{\baselinestretch}{1.05}
\setlength\parindent{2mm}
\setlength{\footskip}{40pt} %set footer size bigger to avoid overlap with text
\usepackage{geometry}
\usepackage{multicol}
\geometry{
    a4paper,
    left=15mm,
    right=15mm,
    top=15mm,
    bottom=15mm,
}

% ----------- Table of Contents -----------
\usepackage{titletoc}
\setcounter{tocdepth}{2} % include subsections in TOC
%
 \titlecontents{section}
 [0] % i
 {\vspace{0.05cm}}
 {\thecontentslabel\enspace}%\thecontentslabel
 {}
 {\titlerule*[0.1cm]{.}\contentspage}%]

 \titlecontents{subsection}
 [1em] %
 {\vspace{0.05cm}}
 {\thecontentslabel\enspace}%\thecontentslabel
 {\hspace*{1em}}
 {\titlerule*[0.1cm]{.}\contentspage}


% ----------- Typography -----------
\usepackage{microtype} % subtle spacing & justification improvements

% ----------- Headings -----------
\usepackage{titlesec}
\titleformat{\section}{\large\bfseries}{\thesection}{1em}{}
\titleformat{\subsection}{\normalsize\bfseries}{\thesubsection}{0.75em}{}

% ----------- Lists -----------
\usepackage{enumitem}
\setlist[itemize]{
    leftmargin = *,
    listparindent = 10mm,
    label = {\listMark} % makes the default item marker a dash
}

% ----------- Figures & Tables -----------
\usepackage{graphicx}
\usepackage{subcaption}
\usepackage{tcolorbox}
\usepackage{float}
\usepackage{tabularx}
\usepackage{booktabs} % professional tables
\usepackage[table]{xcolor}
\usepackage{environ}
\usepackage{tikz}

% Caption formatting
\usepackage{caption}
\captionsetup[table]{
    name=Tabel,                  % rename "Table" to "Tabel"
    labelfont={bf},              % "Tabel 1:" in bold
    textfont={normalfont},       % caption text in normal (non-bold)
    labelsep=colon               % ensures "Tabel 1:" (with colon)
}

\captionsetup[figure]{
    name=Figuur,                 % rename "Figure" to "Figuur"
    labelfont={bf},              % "Figuur 1:" in bold
    textfont={normalfont},       % caption text in normal (non-bold)
    labelsep=colon               % ensures "Figuur 1:" (with colon)
}

% ----------- Footnotes -----------
\usepackage[hang,flushmargin]{footmisc}

% ----------- Hyperlinks -----------
\usepackage[
    colorlinks=true,
    linkcolor=blue,
    citecolor=blue,
    urlcolor=blue,
]{hyperref}

% ----------- Column balance -----------
\usepackage{balance}

% ----------- Bibliography -----------
\usepackage[backend=biber,style=apa]{biblatex}
\DeclareLanguageMapping{dutch}{dutch-apa}
\DefineBibliographyStrings{dutch}{andothers = {et al.}}
\addbibresource{main.bib}
\setlength\bibitemsep{0.25em}

% ----------- ACM-style tweaks (optional) -----------
\settopmatter{printacmref=false}
\setcopyright{none}
\settopmatter{printfolios=true}
\renewcommand\footnotetextcopyrightpermission[1]{}

% ----------- Editors/general version -----------
\usepackage{comment}
\usepackage{ifthen}
\usepackage[utf8]{inputenc}
\usepackage{glossaries}
\newboolean{editorsversion}
\setboolean{editorsversion}{\editorsVersion}  % change to false for student version


% ----------- Paper layout information -----------
\title{\documentTitle}
\subtitle{\documentSubtitle}

\author{\authorName}
\date{\dueDate}
\email{\authorEmail}
\affiliation{
    \institution{\institutionName}
    \city{\institutionCity}
    \country{\institutionCountry}
}

\hypersetup{
    pdfauthor={\authorName},
    pdftitle={\documentTitle},
    pdfborder={0 0 0}
}

% When \begin{document} is written, also add \maketitle and \tableofcontents there. (if settings are true)
\AtBeginDocument
{%
    \ifthenelse{\equal{\UseGlossery}{true}}{%
        \makeglossaries
        \include{src/glossery}
    }{}
    \ifthenelse{\equal{\makeTitlePage}{true}}{%
        \maketitle
    }{}
    \ifthenelse{\equal{\makeTOCpage}{true}}{%
        \tableofcontents
        \newpage
    }{}
}%

% Before \end{document}, add bibliography
\AtEndDocument
{%
    \ifthenelse{\equal{\makeBibliography}{true}}{%
        \printbibliography
    }{}
    \ifthenelse{\equal{\UseGlossery}{true}}{%
        \printglossaries
    }{}
}%
 % import custom general paper settings
% ============================================================
% LATEX CUSTOM COMMANDS LIBRARY
% Author: Matt
% Version: v1.0
% ------------------------------------------------------------
% CONTENTS:
%   1.  SPACING & LAYOUT UTILITIES
%   2.  TEXT PLACEHOLDER & FORMATTING UTILITIES
%   3.  EDITOR TOOLS & COMMENTS
%   4.  TEXT PROCESSING UTILITIES
%   5.  SIMPLE TABLE ENVIRONMENT
%
% This library provides a modular set of reusable LaTeX commands
% to streamline writing, reviewing, and formatting workflow.
%
% ============================================================


% ============================================================
% 1. SPACING & LAYOUT UTILITIES
% ============================================================
%
% These commands simplify vertical spacing and layout adjustments.
% Instead of manually using \vspace repeatedly, these commands
% make your document cleaner and more readable.
%
% Usage examples:
%
%   \vsmall    → Small vertical gap (~0.1 cm)
%   \vmed      → Medium vertical gap (~0.3 cm)
%   \vlarge    → Large vertical gap (~0.6 cm)
%
% ============================================================
\newcommand{\vsmall}{\vspace{0.1cm}}
\newcommand{\vmed}{\vspace{0.3cm}}
\newcommand{\vlarge}{\vspace{0.6cm}}
\newlength{\remainingspace}%
\setlength{\remainingspace}{\dimexpr\pagegoal-\pagetotal-1em\relax}%

% ------------------------------------------------------------
% Centered italicized question or section prompt
%
% Usage:
%   \question{What is the impact of this design choice?}
% ------------------------------------------------------------
\newcommand{\question}[1]{%
    \vsmall%
    \begin{center}%
        \parbox{0.8\linewidth}{\centering\textit{#1}}%
    \end{center}%
    \vsmall%
}


% ============================================================
% 2. TEXT PLACEHOLDER & FORMATTING UTILITIES
% ============================================================
%
% Commands for generating placeholder text, APA-style quotes,
% and custom citation formatting.
%
% Includes:
%   - \lorem                → Generates a block of lorem ipsum text
%   - \uitspraak{txt}{key}{p.}  → APA-style quotes with word count logic
%   - \prepostparencite[...]{}{} → Adds pre/post text to citations
%
% ============================================================

% --- Lorem Ipsum Placeholder --------------------------------
%
% Usage:
%   \lorem
%
% Outputs a single paragraph of lorem ipsum placeholder text.
% ------------------------------------------------------------
\newcommand{\lorem}{%
    Lorem ipsum dolor sit amet, consectetur adipiscing elit. Nullam eget tortor a urna ornare pellentesque.
    Integer sit amet purus nec sem iaculis euismod. Duis at ipsum eu libero pharetra egestas.
    Quisque eleifend odio velit, at sollicitudin metus dictum eu. Integer nec mi congue,
    gravida nibh sed, faucibus mauris. Sed vel ipsum lobortis felis gravida dignissim.
    Curabitur vestibulum turpis eu orci lacinia consectetur.
}

% --- APA Style Quote Command --------------------------------
%
% Automatically formats quotes following APA style.
% Short quotes (<40 words) are inline; long quotes are right-aligned blocks.
%
% Usage:
%   \uitspraak{<quote text>}{<bibkey>}{<page>}
%
% Example:
%   \uitspraak{Design is intelligence made visible.}{rand2020}{42}
% ------------------------------------------------------------
\newcommand{\uitspraak}[3]{%
    \StrCount{#1}{ }[\woordenaantal]%
    \ifthenelse{\woordenaantal < 39}{%
        ``\textit{#1}'' (\citeauthor{#2}, \citeyear{#2}, p.~#3)%
    }{%
        \vsmall%
        \begin{flushright}%
            \parbox{0.95\linewidth}{\textit{#1}}\newline(\citeauthor{#2}, \citeyear{#2}, p.~#3)%
        \end{flushright}%
        \vsmall%
    }%
}

% --- Parenthetical Citation with Pre/Post Text --------------
%
% Usage:
%   \prepostparencite[see][, for more info]{key}
%
% Example:
%   \prepostparencite[see][, for an overview]{doe2022}
%
% Produces: (see Doe, 2022, for an overview)
% ------------------------------------------------------------
\NewDocumentCommand{\prepostparencite}{oom}{%
    (#1 \citeauthor{#3}, \citeyear{#3}#2)%
}


% --- Custom command: Title (Author, Date) ---
% Usage: 
%   \titlectite{key}
%
% Example:
%   \titlectite{doe2022}
%
% Produces: Title (Author, Year)
% ------------------------------------------------------------
\newcommand{\titlectite}[1]{%
    \citetitle{#1} (\citeauthor{#1}, \citefield{#1}{year})%
}

% --- text with lines on the left and right---- --------------
%
% Usage:
%   \ruleline{text in the middle}
%
% Produces: --------- text in the middle ---------
% |
% \> where the line is columnwidth
% ------------------------------------------------------------
\newcommand*\ruleline[1]{\par\noindent\raisebox{.8ex}{\makebox[\linewidth]{\hrulefill\hspace{1ex}\raisebox{-.8ex}{\textbf{#1}}\hspace{1ex}\hrulefill}}}

% ============================================================
% 3. EDITOR TOOLS & COMMENTS
% ============================================================
%
% This block defines tools for editors and reviewers to leave
% inline comments, boxed remarks, or red footnotes.
%
% Controlled by the boolean \boolean{editorsversion}.
% When false, all editor content is hidden.
%
% Provides:
%   - \editorsonly{<text>}       → Inline comment
%   - \begin{editorsonlyBox}...  → Boxed comment block
%   - \editorsfootnote{<text>}   → Red footnote
%
% ============================================================
\ifthenelse{\boolean{editorsversion}}{%
% -------------------- EDITORS VERSION --------------------
    \newcommand{\editorsonly}[1]{%
        \par\vspace{0.05in}%
        {\footnotesize\noindent\textcolor{red!70!black}{\textbf{Editor's note:} #1}}%
        \par\vspace{0.05in}%
    }

    \newenvironment{editorsonlyBox}{%
        \par\medskip%
        \begin{tcolorbox}[%
            enhanced,
            breakable,
            colback=red!5,
            colframe=red!20!white,
            boxrule=0pt,
            borderline north={1pt}{0pt}{red!40!white},
            sharp corners,
            before skip=6pt,
            after skip=6pt,
            title={\textcolor{red!60!black}{\footnotesize Comment for editors:}},
            coltitle=red!70!black,
            fonttitle=\bfseries,
            top=2pt,
            bottom=2pt,
            left=4pt,
            right=4pt
        ]%
        }{%
        \end{tcolorbox}%
    }

    \newcommand{\editorsfootnote}[1]{%
        \textcolor{red!70!black}{\footnote{\textcolor{red!70!black}{Editor's note: #1}}}%
    }

    % Automatically cite all refs in .bib (for review phase)
    \nocite{*}

}{%
% -------------------- STUDENT VERSION --------------------
    \newcommand{\editorsonly}[1]{\ignorespaces}
    \newenvironment{editorsonlyBox}{\comment}{\uncomment}
    \newcommand{\editorsfootnote}[1]{}
}


% ============================================================
% 4. TEXT PROCESSING UTILITIES
% ============================================================
%
% Utilities that manipulate or conditionally display text.
% Includes text splitting, and empty-argument checking.
%
% ============================================================

% --- Split Text Every N Characters --------------------------
%
% Splits input text into lines of maximum <n> characters.
%
% Usage:
%   \spliteveryn{<n>}{<text>}
%
% Example:
%   \spliteveryn{10}{This is an example of splitting text.}
%
% ------------------------------------------------------------
\ExplSyntaxOn
\NewDocumentCommand{\spliteveryn}{mm}{
    \str_set:Nn \l_tmpa_str { #2 }
    \int_zero:N \l_tmpa_int
    \str_map_inline:Nn \l_tmpa_str{
        ##1
        \int_incr:N \l_tmpa_int
        \int_compare:nNnT { \l_tmpa_int } = { #1 }{
            \\
            \int_zero:N \l_tmpa_int
        }
    }
}
\ExplSyntaxOff

% --- Conditional Execution if Non-Empty ---------------------
%
% Usage:
%   \IfNonEmpty{<text>}{<action>}
%
% Executes <action> only if <text> is non-empty.
% Commonly used for optional arguments in macros.
%
% Example:
%   \IfNonEmpty{caption}{\caption{caption}}
%
% ------------------------------------------------------------
\makeatletter
\newcommand{\IfNonEmpty}[2]{%
    \if\relax\detokenize{#1}\relax%
    % empty → do nothing
    \else%
    #2%
    \fi%
}
\makeatother


% ============================================================
% 5. SIMPLE TABLE ENVIRONMENT
% ============================================================
%
% A consistent, minimal table layout system with helper macros
% for headers, subheaders, highlights, and notes.
%
% Usage:
%   \begin{SimpleTable}[<col spec>]{<caption>}{<label>}
%       \TableHeader{Col A & Col B}
%       \TableRow{Data 1 & Data 2}
%       \TableSubHeader{Sub A & Sub B}
%       \TableHighlightRow{Important & Entry}
%       \TableNote{This is a note below the table.}
%   \end{SimpleTable}
%
% Arguments:
%   [<col spec>] : tabularx column layout (default: s{1} s{1})
%   {<caption>}  : optional table caption
%   {<label>}    : optional table label for referencing
%
% ============================================================
\newcolumntype{s}[1]{>{\hsize=#1\hsize\relax\arraybackslash}X}

\ExplSyntaxOn
\NewExpandableDocumentCommand{\GenerateColumns}{m}
{
    \clist_set:Nn \l_tmpa_clist { #1 }
    \clist_use:Nn \l_tmpa_clist { ~s }
% OR simpler:
    % \clist_map_inline:Nn \l_tmpa_clist { s{##1} }
}
\ExplSyntaxOff

\NewEnviron{SimpleTable}[3][s{1} s{1}]{%
    \scalefont{0.75}
    \begin{tabularx}{\linewidth}{#1}
        \BODY
        \arrayrulecolor{black!80}\specialrule{1pt}{0pt}{0pt}\arrayrulecolor{black}
    \end{tabularx}
    \IfNonEmpty{#2}{\captionof{table}{#2}}%
    \IfNonEmpty{#3}{\label{#3}}%
}

% --- Table Styling Macros ------------------------------------------------------

% ========== Header Row ==========
\NewDocumentCommand{\TableHeader}{>{\SplitArgument{7}{&}}m}{%
    \tableheaderaux#1%
}
\NewDocumentCommand{\tableheaderaux}{mmmmmmmm}{%
    \IfValueT{#1}{\textbf{#1}}%
    \IfValueT{#2}{& \textbf{#2}}%
    \IfValueT{#3}{& \textbf{#3}}%
    \IfValueT{#4}{& \textbf{#4}}%
    \IfValueT{#5}{& \textbf{#5}}%
    \IfValueT{#6}{& \textbf{#6}}%
    \IfValueT{#7}{& \textbf{#7}}%
    \IfValueT{#8}{& \textbf{#8}}%
    \\ \arrayrulecolor{black!80}\specialrule{1pt}{0pt}{0pt}\arrayrulecolor{black}
}

% ========== Normal Row ==========
\NewDocumentCommand{\TableRow}{>{\SplitArgument{7}{&}}m}{%
    \tablerowaux#1%
}
\NewDocumentCommand{\tablerowaux}{mmmmmmmm}{%
    \IfValueT{#1}{#1}%
    \IfValueT{#2}{& #2}%
    \IfValueT{#3}{& #3}%
    \IfValueT{#4}{& #4}%
    \IfValueT{#5}{& #5}%
    \IfValueT{#6}{& #6}%
    \IfValueT{#7}{& #7}%
    \IfValueT{#8}{& #8}%
    \\ \arrayrulecolor{black!10}\specialrule{0.5pt}{1pt}{0pt}\arrayrulecolor{black}
}

% ========== Empty Row ==========
\NewDocumentCommand{\TableEmpty}{}{\\}

% ========== Subheader Row ==========
\NewDocumentCommand{\TableSubHeader}{>{\SplitArgument{7}{&}}m}{%
    \tablesubheaderaux#1%
}
\NewDocumentCommand{\tablesubheaderaux}{mmmmmmmm}{%
    \IfValueT{#1}{\underline{\textbf{#1}}}%
    \IfValueT{#2}{& \underline{\textbf{#2}}}%
    \IfValueT{#3}{& \underline{\textbf{#3}}}%
    \IfValueT{#4}{& \underline{\textbf{#4}}}%
    \IfValueT{#5}{& \underline{\textbf{#5}}}%
    \IfValueT{#6}{& \underline{\textbf{#6}}}%
    \IfValueT{#7}{& \underline{\textbf{#7}}}%
    \IfValueT{#8}{& \underline{\textbf{#8}}}%
    \\[1pt]
}

% ========== Highlight Row ==========
\NewDocumentCommand{\TableHighlightRow}{>{\SplitArgument{7}{&}}m}{%
    \tablehighlightaux#1%
}
\NewDocumentCommand{\tablehighlightaux}{mmmmmmmm}{%
    \IfValueT{#1}{\textbf{#1}}%
    \IfValueT{#2}{& \textbf{#2}}%
    \IfValueT{#3}{& \textbf{#3}}%
    \IfValueT{#4}{& \textbf{#4}}%
    \IfValueT{#5}{& \textbf{#5}}%
    \IfValueT{#6}{& \textbf{#6}}%
    \IfValueT{#7}{& \textbf{#7}}%
    \IfValueT{#8}{& \textbf{#8}}%
    \\ \arrayrulecolor{black!50}\specialrule{0.8pt}{1pt}{2pt}\arrayrulecolor{black}
}

% ========== Note Row ==========
\NewDocumentCommand{\TableNote}{m}{%
    \\[-2pt]\multicolumn{99}{l}{\footnotesize\textit{#1}} \\[3pt]%
}

% ========== Section Break ==========
\NewDocumentCommand{\TableSectionBreak}{}{%
    \\[-1pt]\arrayrulecolor{black!70}\specialrule{1.2pt}{3pt}{3pt}\arrayrulecolor{black}
}

% ========== Spacer ==========
\NewDocumentCommand{\TableSpacer}{}{\\[-3pt]}

% ============================================================
% 5. Example enviroment
% ============================================================
%
% A small added enviroment for creating a small example
%
% Usage:
%   \begin{voorbeeld}[<label>]
%       Just the text you want
%   \end{voorbeeld}
%
% Arguments:
%   [label] : The label you want ontop of the voorbeeld box
%
% ============================================================

\newenvironment{voorbeeld}[1]
{
    \vsmall
    \raggedright Voorbeeld \IfNonEmpty{#1}{\textbf{(#1)}}: \\
    \vspace{1mm}
    \hline
    \vspace{1mm}
}
{
    \vspace{1mm}
    \\ \hline
    \vsmall
}
 % import custom commands and packages

\def\documentTitle      {Samenvatting Organisaties en ICT} % replace with actual document title
\def\documentSubtitle   {} % replace with actual subtitle if any, else leave empty
\def\authorName         {Tim v. Naarden} % replace with actual author name
\def\authorEmail        {T.vannaarden@students.uu.nl} % replace with actual author email
\def\institutionName    {Universiteit Utrecht} % replace with actual institution name
\def\institutionCountry {Nederland} % replace with actual institution location
\def\institutionCity    {Utrecht} % replace with actual institution location

\def\UseGlossery        {true} % true om een begrippenlijst aan het einde toe tevoegen 
\def\editorsVersion     {false} % true om editorsnote e.d. te tonen, false voor verbergen
\def\makeTitlePage      {true}  % true om titelpagina te maken, false om over te slaan
\def\makeTOCpage        {true}  % true om inhoudsopgave te maken, false om over te slaan
\def\makeBibliography   {false}  % true om bibliografie te maken, false om over te slaan
\def\listMark           {-} % itemize marker, e.g. '-', '*', '\textbullet', etc.
% -----------------------------------------
\usepackage{float}
\usepackage{listings}

\begin{document}
    \onecolumn
    \section{Lecture 1}
    \subsection{What are Organizations}
        Organizations are social entity's, it consists of people and their relationships.
        Organizations are designed a structured/coordinated activity system which is goal directed 
        and linked to the external environment thru partners, customers etc. 

        Not all organizations are the same, they can differ in Size, product, sector etc. 
        
        Organizations are systems. A system is a regularly interacting or interdependent group of 
        items forming a unified whole whole surrounded and influenced by an environment.
        
        There are two systems perspectives. 
        \begin{enumerate}
            \item{Closed systems perspective: Focus only on the inside of the organization. It is self-contained 
            and sealed of to the outside. (Old fashioned view of organizations)}
            \item{Open systems perspective: Organizations have inputs and outputs, which the organisational design 
            shall manage.}
        \end{enumerate}

        A organization can have multiple business objectives:
        \begin{enumerate}
            \item{Operantional excellence}
            \item{New products, services and business models}
            \item{Customer and supllier intimacy}
            \item{Imrpoved decision making}
            \item{Competitive advantage}
            \item{Survival}
            \item{Becoming more responsible}
        \end{enumerate}
        They also face some challenges:
        \begin{enumerate}
            \item{Globalisation}
            \item{Ethics and social responibility}
            \item{Responsiveness to changes}
            \item{The digital workspace}
            \item{Diversity}
        \end{enumerate}
        Organizations exist to bring resources together(Like building an aircraft carrie) or 
        boosting innovation by delivering new goods \& services. 

        Organizations have become so important because the make use of modern technologies and 
        create value for owners, customers and employees.
    
\end{document}